\chapter{Spin states, auxiliary fermions, and projection}
\label{narr:physical}

A ground-state energy is a statement about a Hamiltonian \emph{and} its state
space. Enlarging the state space may simplify a calculation while changing the
minimum being computed. We first make the enlargement explicit, so that a later
Gaussian approximation cannot be mistaken for an exact solution of the spin
problem.

% EDITORIAL QUESTION: What is a physical trial state and why does it give an upper bound?
\section{The physical variational problem}
The spin Hilbert space is $\mathcal H_s=(\CC^2)^{\otimes N}$. A state is a positive
operator $\rho_s$ of trace one. For a self-adjoint $H$ on this space,
\begin{equation}
 E_0=\min\spec H,\qquad \Tr(\rho_s H)\geq E_0.
 \label{eq:narr:variational}
\end{equation}
The inequality follows by expanding in an energy eigenbasis: the expectation
is a convex combination of eigenvalues. It is exact for any trial density
operator on this Hilbert space. It does not require Gaussianity, convergence of
an iteration, or a small residual.

Restricting the trial states can only increase the infimum, provided they remain
physical states of the same Hamiltonian. An enlarged-space density operator
need not satisfy that proviso. Consequently, the phrase ``variational mean-field
energy'' must specify the space over which variation was performed.

% EDITORIAL QUESTION: How does a four-Majorana enlargement recover a spin algebra?
\section{A local representation and its physical subspace}
% SOURCE [VERIFIED]: Kitaev, cond-mat/0506438v3, Majorana representation; local signs derived in Appendix finite foundations.
Introduce four Hermitian Majoranas per site, ordered locally as
$(b_i^x,b_i^y,b_i^z,c_i)$, satisfying the global canonical anticommutation
relations $\{\gamma_a,\gamma_b\}=2\delta_{ab}\Id$. Define
\begin{equation}
 D_i=b_i^x b_i^y b_i^z c_i,\qquad
 \widetilde\sigma_i^\alpha=\ii b_i^\alpha c_i,\qquad
 P_G=\prod_i\frac{\Id+D_i}{2}.
 \label{eq:narr:projection}
\end{equation}
The $D_i$ are commuting Hermitian involutions. The auxiliary Fock space has
dimension $4^N$; $P_G$ has rank $2^N$. Each spin representative commutes with
all $D_j$, and direct multiplication gives
\begin{equation}
 \widetilde\sigma_i^x\widetilde\sigma_i^y
   =\ii D_i\widetilde\sigma_i^z.
 \label{eq:narr:pauli}
\end{equation}
Thus the usual Pauli algebra holds on $\Ran P_G$. The auxiliary representation
contains redundant states; Eq.~\eqref{eq:narr:pauli} explains exactly which
restriction restores the physical local algebra.

The sign relating $D_i$ to fermion parity depends on the declared ordering and
pairing. With $f_1=(b^x+\ii b^y)/2$ and $f_2=(b^z+\ii c)/2$ at one site,
$D=-(1-2n_1)(1-2n_2)$. In this convention the physical local sector has odd
occupation parity. This statement must not be transferred to another basis by
inspection. On a fixed gauge background, matter-parity selection also depends
on the gauge, labeling, and boundary conditions; the primary-source parity
formula requires those data~\cite{Pedrocchi2011}.

% EDITORIAL QUESTION: Does enforcing an average constraint make an auxiliary state physical?
\section{Projection and the meaning of an energy bound}
For an auxiliary density operator with
$Z_G=\Tr(P_G\rho_{\rm aux})>0$, the normalized projected state is
\begin{equation}
 \rho_G=\frac{P_G\rho_{\rm aux}P_G}{Z_G}.
 \label{eq:narr:normalized_projection}
\end{equation}
An identified spin representation on $\Ran P_G$ makes its energy a physical
variational upper bound. Both the normalization and the condition $Z_G>0$
are necessary. A zero projection cannot define a trial state.

Average constraints generally give less information than support in a joint
kernel. For example, an observable with eigenvalues $+1$ and $-1$ has zero mean
in their equal mixture, though that mixture is not supported in its zero
eigenspace. There is a useful sharper fact for an involution: $\langle D\rangle=1$
\emph{does} force support in its $+1$ eigenspace, because $\Id-D$ is positive.
One must identify which constraint the algorithm actually imposes, rather than
using these different statements interchangeably.

Projection can destroy Gaussianity. At one site the maximally mixed auxiliary
state has zero covariance and obeys Wick's rule. Project it onto $D=+1$ using
the ordering above. The resulting equal mixture of $|01\rangle$ and
$|10\rangle$ still has zero covariance, but its four-Majorana expectation is
$\langle D\rangle=1$, while a Wick contraction of that covariance is zero.
A projected calculation therefore requires more than reusing an unprojected
covariance in the same Wick formulas. Appendix~\ref{app:finite_foundations}
gives the explicit construction and the variational proof.
